\documentclass[handout]{beamer}
\mode<presentation>
\usetheme[secheader]{Boadilla}
\usecolortheme{whale}

\usepackage[utf8]{inputenc}
\usepackage[T1]{fontenc}
\usepackage[indonesian]{babel}
\usepackage{lmodern}
\usepackage{graphicx}
\usepackage{bookmark}

%#make sure to change this part, since it is predefined
%\defbeamertemplate*{footline}{infolines theme}
\setbeamertemplate{footline}
{
	\leavevmode%
	\hbox{%
		\begin{beamercolorbox}[wd=.333333\paperwidth,ht=2.25ex,dp=1ex,center]{author in head/foot}%
			\usebeamerfont{author in head/foot}\insertshortauthor%~~(\insertshortinstitute)
		\end{beamercolorbox}%
		\begin{beamercolorbox}[wd=.433333\paperwidth,ht=2.25ex,dp=1ex,center]{title in head/foot}%
			\usebeamerfont{title in head/foot}\insertshorttitle
		\end{beamercolorbox}%
		\begin{beamercolorbox}[wd=.233333\paperwidth,ht=2.25ex,dp=1ex,right]{date in head/foot}%
			\usebeamerfont{date in head/foot}\insertshortdate{}\hspace*{2em} \insertframenumber{} / \inserttotalframenumber\hspace*{2ex} 
	\end{beamercolorbox}}%
	\vskip0pt%
}


\title{Pemrograman Web}
\subtitle{Kontrak Kuliah}
\author{Cecep Suwanda, S.Si., M.Kom.}
\institute{Universitas Bale Bandung}
\date{\today}

\begin{document}
	
	% Slide Judul
	\begin{frame}
		\titlepage
	\end{frame}
	
	% Slide Informasi Mata Kuliah
	\begin{frame}{Informasi Mata Kuliah}
		\begin{itemize}
			\item \textbf{Kode:} TIF-XXX
			\item \textbf{Semester:} 3 (Tiga)
			\item \textbf{SKS:} 2 SKS
			\item \textbf{Jumlah pertemuan:} 16 pertemuan
			\item \textbf{Komposisi:} 70\% Praktik, 30\% Teori
		\end{itemize}
	\end{frame}
	
	% Slide Informasi Pribadi
	\begin{frame}{Informasi Pribadi}
		
		\begin{itemize}
			\item \textbf{Nama:} Cecep Suwanda, S.Si, M.Kom.
			\item \textbf{Tempat, Tanggal Lahir:} Bandung, 2 Oktober 1982
			\item \textbf{Alamat:} Jl. Radio Babakan Tanjung No. 28 Rt. 01 Rw. 11, Palasari
			\item \textbf{Email:} cecepsuwanda@yahoo.com
			\item \textbf{Telepon:} +62 81320400720
		\end{itemize}
	\end{frame}
	
	% Slide Pendidikan
	\begin{frame}{Pendidikan}
		\begin{itemize}
			\item \textbf{S2:} Teknik Informatika, Fakultas Pasca Sarjana, UNLA, 2024
			\item \textbf{S1:} Matematika, FMIPA UNPAD, 2006
		\end{itemize}
	\end{frame}
	
	% Slide Pengalaman Mengajar
	\begin{frame}{Pengalaman Mengajar}
		\begin{itemize}
			\item Dosen di Teknik Informatika FTI, Universitas Bale Bandung (2025 -- Sekarang)
			\item Dosen di Matematika FMIPA, Universitas Bale Bandung (2008 -- 2023)
		\end{itemize}
	\end{frame}

	% Slide Kontak
	\begin{frame}{Kontak}
		\begin{itemize}
			\item \textbf{Email:} cecepsuwanda@yahoo.com
			\item \textbf{Telepon:} +62 81320400720
			\item \textbf{Alamat:} Jl. Radio Babakan Tanjung No. 28 Rt. 01 Rw. 11, Palasari
		\end{itemize}
	\end{frame}
	
	% Slide Kontrak Perkuliahan (Penilaian dan Aturan Kuliah)
	\begin{frame}{Kontrak Perkuliahan}
		\begin{block}{Penilaian}
			\begin{itemize}
				\item \textbf{UAS:} 25\% dari nilai akhir.
				\item \textbf{UTS:} 20\% dari nilai akhir.
				\item \textbf{Proyek:} 20\% dari nilai akhir.
				\item \textbf{Partisipasi:} 15\% dari nilai akhir.  (edLink)
				\item \textbf{Quiz:} 10\% dari nilai akhir.  (edLink)(per-pertemuan)
			    \item \textbf{Tugas:} 10\% dari nilai akhir. (edLink) 
			    
			\end{itemize}
		\end{block}
		\end{frame}
	
	\begin{frame}{Kontrak Perkuliahan}		
		\begin{block}{Aturan Kuliah}
			\begin{itemize}
				\item Memilih seorang Ketua Kelas.
				\item Ketua Kelas membuat group WA perkuliahan, nama group ``Pemrograman Web 2025''. 
				\item Informasi perkuliahan disampaikan melalui group WA.
				\item Ketua Kelas menghubungi dosen H-1 perkuliahan untuk konfirmasi kuliah atau tidak.
				\item Selalu cek edLink.
				\item Kuliah tatap muka dilaksanakan tepat waktu, dengan toleransi keterlambatan 15 menit, lewat 15 menit jangan masuk kelas.
					\item Tidak ada perbaikan nilai setelah nilai akhir diumumkan.							
			\end{itemize}
		\end{block}
	\end{frame}

	\begin{frame}{Proyek}
	\begin{itemize}
		\item Membuat aplikasi sistem informasi akademik (SIA) sederhana berbasis web.
		\item Boleh menggunakan IDE AI.
		\item Progress proyek diupload ke github.
		\item Membuat dokumentasi proyek (perancangan database, perancangan UI/UX, security, testing) dalam bentuk PDF .
		\item Batas akhir proyek ini 1 minggu sebelum UAS.
	\end{itemize}
\end{frame}

	\begin{frame}{Materi Kuliah}
		\begin{itemize}
			\item \textbf{Pengantar \& Perancangan:} Arsitektur client-server, kebutuhan pengguna, wireframe, mockup, aksesibilitas (a11y), prinsip UX
			\item \textbf{HTML5:} Struktur dokumen, elemen semantik, form, media
			\item \textbf{CSS3:} Selector, box model, Flexbox, Grid, desain responsif
			\item \textbf{JavaScript:} DOM, event handling, validasi form; library/framework front-end (jQuery/Bootstrap/Vue/React)
			\item \textbf{PHP:} Dasar, pemrosesan form, OOP, session/cookie, validasi input
			\item \textbf{MySQL:} CRUD, koneksi PHP--MySQL (MySQLi/PDO)
			\item \textbf{Integrasi \& Deployment:} API, pengujian, keamanan (XSS/CSRF), deployment, administrasi web
		\end{itemize}
	\end{frame}
	
\end{document}
